\documentclass[12pt]{article}
\usepackage[utf8]{inputenc}
\usepackage[T1]{fontenc}
\usepackage{hyperref}
\title{Pidgin Bantu Language Resources for K-12 Curriculum Development}
\author{Compiled by Claude Sonnet 5.5 (claude-sonnet-5-5) for Andr\'as Kornai}
\date{}
\begin{document}
\maketitle

\section{Introduction}
Pidgin Bantu (ISO 639-3 \texttt{fng}; Glottocode fana1235), better known as Fanakalo or Fanagalo, is a Zulu-based pidgin with English and a little Afrikaans input, used mainly as a work and trade lingua franca in South African mining. It has no ethnic community and is used as a second language. Search results describe it as roughly 70\% Zulu, 25\% English and 5\% Afrikaans in its mine variety, and scholars disagree on how ``Bantu'' it really is. No school use was found.

\section{Primary Educational Resources}
No primers, textbooks or school programmes for Fanakalo were found in the search results.

\section{Linguistic and Reference Materials}
\begin{itemize}
\item Survey chapter on Fanakalo in the Atlas of Pidgin and Creole Language Structures: \url{https://apics-online.info/surveys/61}
\item Mesthrie, ``How non-Indo-European is Fanakalo pidgin?'' (John Benjamins): \url{https://benjamins.com/catalog/cll.47.05mes}
\item Same chapter, De Gruyter Brill listing: \url{https://www.degruyterbrill.com/document/doi/10.1075/cll.47.05mes/html}
\item The Oxford Guide to the Bantu Languages, chapter discussing Fanakalo: \url{https://academic.oup.com/book/61440/chapter/534769303}
\end{itemize}
A phrase book, grammar and dictionary by J.D. Bold (Johannesburg) is cited in Mesthrie's bibliography according to the search summary, but no page for it was found, so it is not listed as a resource.

\section{Religious and Community Resources}
None found.

\section{Digital and Audio Resources}
None found.

\section{Curriculum Development Context}
Fanakalo has no native-speaker community and is not known to be taught in schools, so a native-language K-12 curriculum is not applicable in the usual sense. Available material is scholarly description only; any classroom use would need new materials, probably built on Zulu resources.

\end{document}